\subsection{THETA3-C32: an explicit exceptional cycle for the accessible cyclic-four Jacobian}\label{C32}

The C25 target admits a constructive higher-correspondence
route, although its C27 map quartic vanishes on the covers
in C28 and C30.  The construction below uses the actual
curve equation, a plane in a symmetric power of its quotient
line, and three blowups.  It does not insert a previously
supplied Weil cycle or use Markman's theorem.

\paragraph{Statement and relation to the published construction.}
Let
\[
 H:\quad y^2=z\prod_{j=1}^4(z^2-a_j^2),
 \qquad a_j\ne0,\quad a_j^2\ne a_k^2\ (j\ne k),
                                                        \tag{C32a}
\]
be defined over $\overline\Q$, and let
$A=J(H)$ with its action by
$\rho(z,y)=(-z,iy)$.  This includes the specific
simple fourfold of C25.  There is an explicitly
constructed integral algebraic cycle
$C_A\in\CH^2(A)$, allowed to have signed components,
such that
\[
 0\ne\operatorname{cl}(C_A)\in W_{\Q(i)}(A),\qquad
                   \int_A\operatorname{cl}(C_A)^2=73728.
                                                        \tag{C32b}
\]
Its pullback by the endomorphism $[2+i]$ supplies an
independent rational direction of the Weil space:
the two eigenvalues there are $(2+i)^4=-7+24i$
and its conjugate.  The scalar automorphism $[i]$
itself acts trivially on this fourth exterior power
and is not used for that step.

The underlying method is Schoen,
\emph{Hodge classes on self-products of a variety
with an automorphism}, Compositio~65 (1988),
Theorem~2.0, printed p.~11, and Corollary~3.1,
p.~24.  Its paired branch hypothesis applies here.
We give a complete specialized construction and
nonvanishing proof.  In particular the exceptional
cohomology argument below treats inertia of order
two inside the cyclic group of order four directly.
It does not assume that such inertia is detected
by a prime-degree quotient.

If $f:X\to H$ is the actual surjection from the
compact arithmetic ball cover in C25, the same
equations produce a surface cycle on $X^4$ whose
image in $A$ is a positive integer multiple of
$C_A$.  The precise multiplier is given below.
This is a higher correspondence on a product of
ball surfaces; the cycle's surface components are
not asserted to be ball quotients.

\paragraph{The cyclic cover and its two-dimensional tensor.}
Put $t=z^2$ and $b_j=a_j^2$.  Equation (C32a)
becomes the connected cyclic-four cover
\[
       y^4=t\prod_{j=1}^4(t-b_j)^2
                                                        \tag{C32c}
\]
of $\mathbf P^1$.  Its six branch points have
weights
\[
                 (1,3,2,2,2,2)\pmod4
                                                        \tag{C32d}
\]
at $0,\infty,b_1,b_2,b_3,b_4$, respectively.
Pair them as $(0,\infty),(b_1,b_2),(b_3,b_4)$.
Each pair has total weight four.  The involution
$\rho^2$ is hyperelliptic, hence acts as $-1$ on
$H^1(H,\Q)$.  Thus the only eigenvalues of $\rho$
on first cohomology are $i,-i$, each with multiplicity
four.  On holomorphic differentials the basis
$z^jdz/y$, $0\leq j\leq3$, gives multiplicity two
for each eigenvalue.  The determinant lines are
therefore of type $(2,2)$.

Let
\[
 N=\{(e_1,e_2,e_3,e_4)\in(\Z/4\Z)^4:
                                      \sum e_i=0\},
 \qquad G=N\rtimes\mathfrak S_4,
 \qquad W_0=H^4/G.
                                                        \tag{C32e}
\]
Write $q:H^4\to W_0$ for the quotient.  Its generic
degree is $|G|=4^3\cdot4!=1536$.  The residual
cyclic group is generated by $\delta$, induced by
applying $\rho$ on the first factor.  The map
\[
             \beta_0:W_0\longrightarrow
                  \operatorname{Sym}^4(\mathbf P^1)=\mathbf P^4
\]
has degree four.  On the open set of four distinct
unramified points its root coordinate is the product
of their four $y$-coordinates; its fourth power is
the product of the four right sides in (C32c).
These function fields determine the normal cover
everywhere.

Rational cohomology of a finite quotient is the
invariant cohomology, as follows by equivariant
triangulation and averaging over the finite group.
The part of $H^4(W_0,\Q)$ on which $\delta^2=-1$
is rationally two-dimensional.  To verify this
without a dimension assumption, use K\"unneth on
$H^4$.  Invariance under $N$ forces the four
cyclic characters in a tensor to agree.  If a
factor has degree zero or two its character is
trivial; such a tensor cannot have primitive
residual character.  The primitive part is thus
in $H^1(H)^{\otimes4}$.  Permutations of the four
odd-degree factors include their Koszul signs,
so their invariant tensors are
\[
     \bigwedge^4 H^1(H,\C)_i
          \ \oplus\ 
     \bigwedge^4 H^1(H,\C)_{-i}.
                                                        \tag{C32f}
\]
Each summand has dimension one.  Denote this
rational two-dimensional subspace by $U_0$.

\paragraph{The explicit plane.}
For each branch point $c$ let $S_c\subset\mathbf P^4$
be the hyperplane of effective degree-four divisors
containing $c$, with defining evaluation form $\ell_c$.
Any five of the six evaluation forms are independent:
a nonzero binary quartic cannot vanish at five
distinct points.  Consequently their unique linear
relation has all coefficients nonzero.  Rescale
them, over $\overline\Q$, to write this relation as
\[
       u_1+v_1+u_2+v_2+u_3+v_3=0,
                                                        \tag{C32g}
\]
where $(u_i,v_i)$ correspond to the three paired
branch points.  This realizes $\mathbf P^4$ as
the indicated hyperplane in $\mathbf P^5$.

Let $P_0\simeq\mathbf P^2$ be the plane
\[
                         u_i+v_i=0\quad(1\leq i\leq3).
                                                        \tag{C32h}
\]
Equivalently it is spanned by the three quartics
whose zero divisors consist of all branch points
except one chosen pair.  On $P_0$, the two
hyperplanes of a pair restrict to the same line;
the three lines are its coordinate lines.
The fourth-root equation restricts to a fourth
power up to a nonzero algebraic scalar, because
the exponents in each pair sum to four.

More explicitly, if $F(T)$ represents a quartic
with leading coefficient $\ell_\infty$, the norm
of the right side of (C32c) is
\[
 \frac{\ell_0\prod_{j=1}^4\ell_{b_j}^2}
                         {\ell_\infty^9}
 =\frac{\ell_0\ell_\infty^3
                 \prod_{j=1}^4\ell_{b_j}^2}
                         {\ell_\infty^{12}}.
                                                        \tag{C32i}
\]
Thus $W_0$ is the normalization of the cyclic
cover with root line $\mathcal O_{\mathbf P^4}(3)$
and branch weights (C32d).  On $P_0$ the numerator
in the last expression is a scalar times
$(u_1u_2u_3)^4$.  A fourth root of that scalar
is the only additional constant needed to label
the four sheets over the plane.

\paragraph{The three blowups and the normal bundle.}
Blow up, in order, the strict transforms of
\[
                 K_i=\{u_i=v_i=0\},\qquad i=1,2,3,
\]
and call the successive smooth varieties $M_i$.
The six hyperplanes and their strict transforms
are normal-crossing divisors; the centers are
intersections of two of them.  Let $P_i$ be the
strict transform of $P_0$.  At the $i$th center
its intersection with $P_{i-1}$ is a Cartier
line, so $P_i\simeq P_0$.  The two strict branch
hyperplanes of that pair no longer meet $P_i$.
After all three steps $P=P_3$ misses all the
strict branch hyperplanes.

Here is also the full local normal-bundle
calculation.  After the first two blowups,
$M_2$ is the incidence variety resolving
\[
        ([u_1:v_1],[u_2:v_2]):\mathbf P^4
                                         \dashrightarrow\mathbf P^1\times\mathbf P^1.
\]
It is a $\mathbf P^2$ bundle: after writing
$(u_i,v_i)=(a_i t_i,b_i t_i)$ for $i=1,2$,
the single relation is
$(a_1+b_1)t_1+(a_2+b_2)t_2+u_3+v_3=0$.
The coefficient of $v_3$ is one, so the kernel
is a rank-three vector bundle.  These incidence
equations are the two blowup charts: pulling
back $(u_1,v_1)$ and then the strict ideal
$(u_2,v_2)$ makes each principal.  The fibre at
$[a_i:b_i]=[1:-1]$ is exactly $P_2$.

Near that fibre use base coordinates $r_1,r_2$
and fibre coordinates $[t_1:t_2:u_3]$, so that
$v_3=-u_3-r_1t_1-r_2t_2$.  The final center has
ideal $(u_3,r_1t_1+r_2t_2)$.  The blowup chart
containing the whole strict plane has equation
\[
                    r_1t_1+r_2t_2=u_3w,
\]
and the plane is $r_1=r_2=w=0$.  Its normal
bundle is therefore the kernel in
\[
 0\longrightarrow N_{P/M_3}\longrightarrow
 \mathcal O_P^3
 \xrightarrow{(t_1,t_2,-u_3)}\mathcal O_P(1)
 \longrightarrow0.
                                                        \tag{C32j}
\]
This is the Euler sequence, so
$N_{P/M_3}\simeq\Omega^1_{\mathbf P^2}(1)$ and
\[
 c(N_{P/M_3})=(1+H)^{-1}=1-H+H^2,
                    \int_P c_2(N_{P/M_3})=1.
                                                        \tag{C32k}
\]

Let $W_i$ be the normalization of the main cyclic
cover after each blowup.  The exceptional divisor
has total branch multiplicity four, so division
of the root by its local exceptional equation
removes this multiplicity.  The new branch
divisor is precisely the union of the strict
transforms, with the original weights.  Hence
$W_3\to M_3$ is \'etale over a neighborhood of
$P$.  Since $\mathbf P^2$ has no nontrivial
connected finite \'etale covers, its preimage
is a disjoint union
\[
                      Q_0,Q_1,Q_2,Q_3\simeq\mathbf P^2,
\]
labelled so that $\delta Q_j=Q_{j+1}$.
The normal bundle of each is (C32j).

\paragraph{The inertia argument, including weights two modulo four.}
We verify that the primitive part of cohomology is
unchanged by every map $\varphi_i:W_i\to W_{i-1}$.
Let $E_i\to K_i$ be the exceptional $\mathbf P^1$
bundle on the base, let $F_i$ be its reduced
inverse image in $W_i$, and let $K_i'$ be its
image in $W_{i-1}$.  Write the paired weights
as $a,4-a$, and put $d=\gcd(4,a)$.  Thus
$a=1$ or $2$ in this construction.

The bundle $E_i$ is the projectivization of the
sum of the two normal line bundles of the
intersecting branch divisors.  Its two sections
come from the strict branch divisors.  All other
branch components restrict to vertical divisors
pulled back from $K_i$, because the original
divisors are normal crossing.  In a ratio
coordinate $t$ on this bundle, the normal cyclic
cover $F_i$ has equation
\[
                       v^4=t^a u,
                                                        \tag{C32l}
\]
where $u$ comes from the base, followed by
normalization.  This description includes points
on the vertical branch divisors.  It also identifies
this normalization with the actual reduced inverse
image of $E_i$: locally analytically one extracts
a fourth root of the unit in the branch equation,
after which the radial exceptional coordinate is
absent from the normalized monomial equation.
Setting that coordinate to zero gives exactly
the displayed normal cyclic cover.

There is a global connected $\mathbb G_m$ action
on $F_i$, fixing $K_i$, with local formula
\[
              t\longmapsto\lambda^{4/d}t,
              \qquad v\longmapsto\lambda^{a/d}v.
                                                        \tag{C32m}
\]
To check gluing, use the diagonal scaling action
on the sum of the two normal line bundles.
The root line on $E_i$ is $\mathcal O_{E_i}(1)$
times a line from $K_i$: before restriction it
is the old root line tensored with
$\mathcal O(-E_i)$.  Give this root line the
linearization of weight $a/d$ in the displayed
chart.  On the other chart $t'=t^{-1}$ the
root is $v'=v/t$, with weight $-(4-a)/d$,
so the formulas agree on the overlap.  Base
transition functions have weight zero.
Normalization is functorial for these
automorphisms and preserves the action.

For $\lambda\in\mu_{4/d}$, the base coordinate
$t$ is unchanged, and the resulting automorphisms
are exactly the deck subgroup
$I=\mu_{4/d}\subset\mu_4$.  Here
$\gcd(a/d,4/d)=1$.  Consequently $I$ acts
trivially on rational cohomology of $F_i$:
each element belongs to a connected group
action and is homotopic to the identity.
It also fixes $K_i'$ pointwise.  Indeed a
$\mathbb G_m$ orbit in $F_i$ maps into one
finite fibre of $K_i'\to K_i$; its image is
connected and is therefore a point.  For
elements of $I$ this is the given equivariant
deck action.  A primitive order-four character
is nontrivial on $I$, both when $I=\mu_4$
and when $I=\mu_2$.  Hence neither
$H^*(F_i,\Q)$ nor $H^*(K_i',\Q)$ has a
primitive component.

The map is an isomorphism off these exceptional
sets, and the quotient pairs give a homeomorphism
$W_i/F_i\simeq W_{i-1}/K_i'$.  This follows
directly from properness and the compact
Hausdorff topology after collapsing each entire
closed exceptional set to one point.  Their
relative cohomologies therefore agree.  Apply
the long exact sequences of the two pairs and
the primitive idempotent.  The preceding
vanishing proves
\[
 \varphi_i^*:H^*(W_{i-1},\Q)_{\rm prim}
                    \xrightarrow{\ \sim\ }
                    H^*(W_i,\Q)_{\rm prim}.
                                                        \tag{C32n}
\]
This supplies the exceptional-cohomology step
without a prime-quotient assumption.  In
particular a weight-two branch of the
cyclic-four cover is not mistakenly declared
visible in its degree-two quotient.

\paragraph{Nonzero cycle and exact intersection number.}
The normal covers $W_i$ have finite quotient
singularities.  Locally, after extracting a
root of the analytic unit, their equations
are fourth roots of products of normal-crossing
coordinates.  Substituting fourth roots for
those coordinates realizes each local branch
as a finite abelian quotient of a polydisc.
They are therefore oriented rational homology
manifolds: rational cohomology of a local
quotient sphere is the invariant cohomology
of the sphere, and a complex linear action
preserves its orientation.  In particular
rational Poincar\'e duality, the degree formula
and the projection formula apply to these
proper maps.  The planes $Q_j$ lie in their
smooth loci, so their ordinary fundamental
classes have their usual self-intersection
interpretation.

Set
\[
              C_3=Q_0-Q_2,\qquad c_3=\operatorname{cl}(C_3).
\]
Its class is primitive, since $\delta^2$ sends
it to its negative.  The two planes are
disjoint, and (C32k) gives
\[
                         \int_{W_3}c_3^2=2.
                                                        \tag{C32o}
\]
Write $\varphi:W_3\to W_0$, put
$C_0=\varphi_*C_3$, and let $c_0$ be its
cohomology class.  Proper cycle classes and
Gysin agree.  The degree-one identity
$\varphi_*\varphi^*=1$, together with (C32n),
implies
\[
       \varphi^*c_0=c_3,\qquad
       c_0\in U_0,\qquad
                         \int_{W_0}c_0^2=2.
                                                        \tag{C32p}
\]
Thus no nonzero primitive class is lost in
blowing the construction down.

To obtain an actual cycle on the smooth
$H^4$, take a smooth projective resolution
$T$ of the main component of
$H^4\times_{W_0}W_3$.  Denote its maps by
$p:T\to H^4$ and $q':T\to W_3$.
The first is birational and
$q\circ p=\varphi\circ q'$.
Each $Q_j\hookrightarrow W_3$ is a regular
codimension-two embedding into a smooth
neighborhood.  Its refined pullback along
$q'$ is consequently an integral cycle on
$T$, including any excess-intersection
terms.  Define
\[
 C_H=p_*\bigl(q'^![Q_0]-q'^![Q_2]\bigr)\in\CH^2(H^4).
                                                        \tag{C32q}
\]
The notation means the refined intersection
with $[T]$, pushed from the fibre product
to $T$.  Naturality of the Thom class for
a regular embedding proves the cycle-class
identity for this operation.  Since
$p_*p^*=1$, the commuting diagram and
(C32p) give
\[
 \operatorname{cl}(C_H)=p_*q'^*c_3
                      =p_*p^*q^*c_0=q^*c_0,
 \qquad
              \int_{H^4}\operatorname{cl}(C_H)^2=3072.
                                                        \tag{C32r}
\]
The last integer is $1536\cdot2$.
This construction does not presume that the
finite quotient map $q$ is flat at its
singular points.

\paragraph{The actual exceptional cycle on the Jacobian.}
Choose a base point on $H$ and let
$a:H^4\to A$ be the sum of the four Abel
maps.  It has generic degree $4!=24$:
$\operatorname{Sym}^4H\to\operatorname{Pic}^4H$
is birational, because a general divisor of
degree $g=4$ has one section by Riemann--Roch.
One can check the latter by choosing four
points which impose independent conditions
on the four-dimensional canonical space.

Pullback by $a$ identifies $W_{\Q(i)}(A)$
with $q^*U_0$.  Indeed it sends a first
cohomology class to the sum of its copies
in the four factors.  The cup product of
two classes with the same eigenvalue $i$
or $-i$ vanishes on one curve, since its
eigenvalue would be $-1$ in the trivial
representation $H^2(H)$.  Expanding a
fourfold wedge therefore leaves precisely
the alternating tensors in (C32f), with
nonzero coefficient.  This is an isomorphism
of the rational two-dimensional spaces.

Define $C_A=a_*C_H$.  If
$\operatorname{cl}(C_H)=a^*w_0$, then
$\operatorname{cl}(C_A)=24w_0$ is a nonzero
pure Weil class.  Finally
\[
 \int_A\operatorname{cl}(C_A)^2
   =24\int_{H^4}\operatorname{cl}(C_H)^2
   =24\cdot3072=73728.
                                                        \tag{C32s}
\]
All maps, planes, blowups, normalizations
and cycles are defined over a finite
extension of the original number field.
The only sheet choice above the plane
is a fourth root of the algebraic scalar
in (C32i); a projective resolution is
available in characteristic zero.

\paragraph{A cycle constructed on the known ball power.}
Let $f:X\to H$ be the actual map from
C25.  It is flat: at a point of the
smooth base curve a uniformizer pulls
back to a nonzerodivisor on the smooth
integral surface, so its local module
is torsion free, hence flat, over the
discrete valuation ring.  Choose a very
ample line bundle $L_X$ and four general
smooth divisors $D_i\in|L_X|$ which
dominate $H$.  Each $D_i\to H$ is finite
flat of the same positive degree
\[
                       d=\deg(L_X|_{X_\eta}).
\]
Pull the explicit cycle (C32q) back to
$D_1\times\cdots\times D_4$, and include
it into $X^4$; denote the resulting
surface cycle by $\Gamma$.  Equivalently
one pulls its equations through the four
ball fibrations and takes one relative
hyperplane section in every ball factor.
The finite degree formula and projection
formula give
\[
       (a\circ f^4)_*\Gamma=d^4 C_A,
 \qquad
       \int_A\operatorname{cl}((a\circ f^4)_*\Gamma)^2
                                      =d^8\cdot73728>0.
                                                        \tag{C32t}
\]
This proves nonvanishing of the actual
higher-correspondence cycle.  If
$L_X=K_X^{\otimes m}$ is chosen very
ample and $e$ is the degree of the
Stein base curve over $H$, C28's genus
$76$ fibre gives the additional explicit
value $d=150me$ by adjunction.

The source of this nonzero class is the
signed split-plane cycle and its refined
pullbacks, not the fundamental class of
a curve product or of a ball surface.
It therefore does not contradict C14,
C28 or C30.  It supplies a concrete
exceptional cycle built from the
accessible curve and ball fibrations,
while the literal single-ball theta
packet comparison remains a separate
problem.

\paragraph{The full C24 family in every even dimension.}
The same calculation gives explicit Weil cycles for
every member of C24, in every dimension $2n\geq4$.
Precisely, let
\[
 H_n:\quad y^2=z\prod_{j=1}^{2n}(z^2-a_j^2),
 \qquad A_n=J(H_n),
                                                        \tag{C32u}
\]
with the same nonzero and distinct-square conditions.
Then an integral codimension-$n$ cycle $C_{A_n}$
is constructed by the preceding operations, and
\[
 \begin{split}
 0\ne\operatorname{cl}(C_{A_n})&\in W_{\Q(i)}(A_n),\\
 \int_{A_n}\operatorname{cl}(C_{A_n})^2
       &=2(-1)^n4^{2n-1}\bigl((2n)!\bigr)^2.
 \end{split}
                                                        \tag{C32v}
\]
The sign is the usual middle-degree primitive
polarization sign; the polarized square is positive.
For the particular non-CM members obtained in C24,
these are actual cycles on the supplied targets.
Their being outside the divisor algebra follows
when that target has the additional N\'eron--Severi
property used in C14.  That property is proved in
C25 for the chosen fourfold; it has not been
proved here for every higher-dimensional C24
member.  Algebraicity and the nonzero pure Weil
projection in (C32v) do not require that extra
assertion.

Here are all changes needed in the proof.  The
cyclic-four equation is
\[
            y^4=t\prod_{j=1}^{2n}(t-a_j^2)^2.
\]
Its branch weights are $1,3$ and $2n$ copies of
$2$, paired into $n+1$ pairs.  Its primitive
first-cohomology rank over $\Q(i)$ is $2n$.
The $2n$ holomorphic differentials $z^jdz/y$
give multiplicities $(n,n)$, so the determinant
lines have type $(n,n)$.

In $\operatorname{Sym}^{2n}\mathbf P^1
=\mathbf P^{2n}$, the $2n+2$ evaluation
hyperplanes have one relation, with every
coefficient nonzero.  After rescaling write
it as $\sum_{i=1}^{n+1}(u_i+v_i)=0$.
The plane $P_0\simeq\mathbf P^n$ is
$u_i+v_i=0$ for all $i$; it is spanned by
the $n+1$ divisors omitting one pair.  The
cyclic root line is $\mathcal O(n+1)$:
the total branch weight is $4(n+1)$.
The same norm calculation as (C32i) shows
that its restriction to $P_0$ is a fourth
power up to an algebraic scalar.

Blow up the $n+1$ paired hyperplane
intersections.  The first $n$ blowups give
the incidence $\mathbf P^n$ bundle over
$(\mathbf P^1)^n$ resolving the first
$n$ pair ratios.  The proof is exactly
the displayed incidence calculation:
\[
       \sum_{i=1}^{n}(a_i+b_i)t_i+u_{n+1}+v_{n+1}=0
\]
is a surjective relation on a rank-$(n+2)$
bundle.  The desired $\mathbf P^n$ is
the fibre where every ratio is $[1:-1]$.
For the final blowup its local equation
near the strict plane is
\[
                  \sum_{i=1}^{n}r_it_i=u_{n+1}w.
\]
Hence its normal bundle is the kernel of
\[
       \mathcal O_{\mathbf P^n}^{n+1}
        \xrightarrow{(t_1,\ldots,t_n,-u_{n+1})}
                          \mathcal O_{\mathbf P^n}(1),
\]
namely $\Omega^1_{\mathbf P^n}(1)$.  Its
top Chern number is $(-1)^n$.
The normalized cyclic cover is \'etale near
the strict plane and splits into four
copies.  Their signed difference $Q_0-Q_2$
has square $2(-1)^n$.

Every exceptional step has weights $1,3$
or $2,2$.  The global action (C32m), the
vanishing of primitive exceptional
cohomology, and the relative-pair proof
apply without any dimension change.
The regular embedding now has codimension
$n$, so the graph-resolution refined
pullback (C32q) works in that codimension.
The finite quotient degree is
$4^{2n-1}(2n)!$, and the addition map
$H_n^{2n}\to J(H_n)$ has degree $(2n)!$.
K\"unneth with the same character and
Koszul-sign argument identifies its
primitive subspace with the Weil space.
Multiplying these two degrees by
$2(-1)^n$ proves (C32v).

The second rational direction can again
be obtained using $[2+i]^*$.  Indeed
$(2+i)^{2n}$ is never real: otherwise
$((2+i)/(2-i))^{2n}=1$, whereas
$(2+i)/(2-i)=(3+4i)/5$ is not a root
of unity in $\Q(i)$, whose roots of
unity are $1,-1,i,-i$.  Thus the two
Weil eigenvalues are distinct and
nonrational, and the two rational
cycle classes are independent.

Finally C24 supplies the actual ball
cover $X_n\to H_n$.  Cutting its
$2n$-fold power by one dominating
very ample curve in each factor and
pulling back the same explicit cycle
gives an $n$-dimensional cycle whose
image in $A_n$ is $d^{2n}C_{A_n}$.
Its square is $d^{4n}$ times (C32v).
Thus the all-dimensional conclusion
is a specialized Schoen construction
on this cyclic-four family, together
with the already constructed ball
access and the exact degree and
intersection normalizations proved
here.  It is not a proof for arbitrary
Weil-type abelian $2n$-folds.

\paragraph{Actual integral specialization at a good prime.}
Let $K$ be a number field over which the chosen curve, its
$\Q(i)$-action and the constructed cycle are defined, and
let $R$ be a discrete valuation ring above a finite prime
of $K$. Suppose $A_n$ extends to an abelian scheme
$\mathcal A_n/R$, with the order action extended. This
holds, for example, at any odd prime where all $a_j$ are
units and all $a_j^2-a_k^2$ are units, after adjoining the
chosen sheet label. Take the schematic closure in
$\mathcal A_n$ of each integral component of $C_{A_n}$,
with its signed multiplicity. Each such closure is flat
over $R$: its local rings are torsion free over the
discrete valuation ring. It has relative dimension $n$.
Its special fibre therefore defines an actual integral
cycle $C_{A_n,s}$ of codimension $n$.

For every $\ell$ different from the residue characteristic,
smooth proper base change and compatibility of cycle
classes with specialization give
\[
 \operatorname{sp}\bigl(\operatorname{cl}_{\ell}(C_{A_n})\bigr)
       =\operatorname{cl}_{\ell}(C_{A_n,s}),\qquad
 \deg\bigl(C_{A_n,s}\cdot C_{A_n,s}\bigr)
       =2(-1)^n4^{2n-1}\bigl((2n)!\bigr)^2.
                                                        \tag{C32w}
\]
Here the intersection is the Chow intersection on the
smooth special-fibre abelian variety; its degree equals
the cup-product trace. These are precisely the
specialization operations and exact SGA locators recorded
in C17. The specialization is equivariant for the extended
order action and the exterior-power decomposition of
$H^1$. Thus the class remains in the rank-two determinant
Weil summand, and (C32w) proves that it is nonzero. The
integer in (C32w) remains a nonzero element of $\Q_\ell$
even when the residue characteristic divides it. No
claim of ordinary proper intersection of the signed
components is needed.

This argument deliberately uses closure of the already
constructed integral cycle; it does not require a new
resolution of a mixed-characteristic quotient. If the
ball surface, its map to the curve and the chosen
dominating divisors also extend smoothly, with the
divisors finite flat of degree $d$, the proper pushforward
and flat pullback identities in (C32t) specialize as
well. Such simultaneous models exist after excluding
finitely many primes, since all the stated maps and
equations have been fixed over a number field. This
gives the same actual ball-derived special-fibre cycle
and the factor $d^{2n}$. Its identification with an
interior Kudla--Rapoport moduli cycle is not a consequence
of specialization alone.

\framework{Schoen's original construction is
Theorem~2.0 and Lemmas~2.3--2.8, printed
pp.~11--19; the passage to Jacobians is
Corollary~3.1, p.~24.  The proof above
specializes the plane and normal bundle
and provides the full order-four inertia
argument.  Its regular-embedding pullback,
proper pushforward and graph resolution
are the correspondence operations in
C17 and C23, with the Grothendieck
anchors cote~44, scans~7--8, and
cote~16, scans~43--46.  No new theorem
about all non-CM Hodge classes is inferred
from this particular cyclic-cover family.}
